% @see : https://coopmaths.fr/alea/?uuid=dnb_2026_01_sujet0va_3&alea=fjpu&uuid=dnb_2024_09_metropole_4&alea=dhH9&uuid=dnb_2022_09_polynesie_4&alea=fEGA&v=eleve&pdfParam=eyJ0aXRsZSI6IiIsInJlZmVyZW5jZSI6IiIsInN1YnRpdGxlIjoiIiwic3R5bGUiOiJDb29wbWF0aHMiLCJmb250T3B0aW9uIjoiU3RhbmRhcmRGb250IiwidGFpbGxlRm9udE9wdGlvbiI6MTIsImR5c1RhaWxsZUZvbnRPcHRpb24iOjE0LCJjb3JyZWN0aW9uT3B0aW9uIjoiQXZlY0NvcnJlY3Rpb24iLCJxcmNvZGVPcHRpb24iOiJTYW5zUXJjb2RlIiwidHlwZUZpY2hlIjoiRmljaGUiLCJkdXJhdGlvbkNhbk9wdGlvbiI6IjkgbWluIiwidGl0bGVPcHRpb24iOiJTYW5zVGl0cmUiLCJuYlZlcnNpb25zIjoxLCJleG9zIjp7fX0\%3D
\documentclass[a4paper,12pt,fleqn]{memoir}
\setlrmarginsandblock{.5cm}{.5cm}{*}
\setulmarginsandblock{.5cm}{.5cm}{*}
 \usepackage{auto-pst-pdf}
%%% COULEURS %%%
\usepackage[table,svgnames]{xcolor}
\definecolor{nombres}{cmyk}{0,.8,.95,0}
\definecolor{gestion}{cmyk}{.75,1,.11,.12}
\definecolor{gestionbis}{cmyk}{.75,1,.11,.12}
\definecolor{grandeurs}{cmyk}{.02,.44,1,0}
\definecolor{geo}{cmyk}{.62,.1,0,0}
\definecolor{algo}{cmyk}{.69,.02,.36,0}
\definecolor{correction}{cmyk}{.63,.23,.93,.06}
% Couleur des filets des tableaux

\usepackage[most]{tcolorbox} % Supprimé à cause de conflits avec ProfCollege
\tcbuselibrary{documentation}

%%%%%%%% Modification paramétrage pour footnotes
% Le paquet semble déjà appelé en amont, curieux qu'il n'y ait pas eu de clash
% avec l'appel ci-dessous mais seulement quand j'ai essayé de l'inclure avec l'option lualatex

% \usepackage{hyperref}
% Parametrages
\hypersetup{
    colorlinks=true,% On active la couleur pour les liens. Couleur par défaut rouge
    linkcolor=black,% On définit la couleur pour les liens internes
    % filecolor=magenta,% On définit la couleur pour les liens vers les fichiers locaux      
    urlcolor=black,% On définit la couleur pour les liens vers des sites web
    % pdftitle={Puissance Quatre},% On définit un titre pour le document pdf
    % pdfpagemode=FullScreen,% On fixe l'affichage par défaut à plein écran
    }
% Pour avoir les numérotations arabic y compris dans les environnements tcolorbox
\renewcommand{\thempfootnote}{\arabic{mpfootnote}}
%%%%%%%%

\usepackage{multicol} 
\usepackage{calc}
\usepackage{enumitem}
\usepackage{graphicx}
\usepackage{tabularx}
%\usepackage{tablvar}

\setlength{\parindent}{0mm}
\renewcommand{\arraystretch}{1.5}
\renewcommand{\labelenumi}{\textbf{\theenumi{}.}}
\renewcommand{\labelenumii}{\textbf{\theenumii{}.}}
\setlength{\fboxsep}{3mm}
 

\setlength{\premulticols}{0pt}

\newcounter{ExoMA}

\newlength{\parindentMA}%
\setlength{\parindentMA}{\parindent} 

\usepackage{simplekv}
 

\usepackage{fancyhdr}
\fancyhead[C]{}
\fancyfoot{}

\def\bla{}
 
 


%%% MISE EN PAGE %%%

\usepackage[margin=1.5cm]{geometry}
\usepackage{setspace}
  \usepackage{eurosym}
\usepackage{pgf}%
\usepackage{pgfplots}
\pgfplotsset{compat=1.18}
\usetikzlibrary{babel,arrows, arrows.meta,calc,fit,patterns,plotmarks,shapes.geometric,shapes.misc,shapes.symbols,shapes.arrows,shapes.callouts, shapes.multipart, shapes.gates.logic.US,shapes.gates.logic.IEC, er, automata,backgrounds,chains,topaths,trees,petri,mindmap,matrix, calendar, folding,fadings,through,positioning,scopes,decorations.fractals,decorations.shapes,decorations.text,decorations.pathmorphing,decorations.pathreplacing,decorations.footprints,decorations.markings,shadows}

\spaceskip=2\fontdimen2\font plus 3\fontdimen3\font minus3\fontdimen4\font\relax %Pour doubler l'espace entre les mots
 
 
  

\usepackage[french]{babel}
% loadPackagesFromContent
\usepackage{pstricks}
\usepackage{pst-plot}
\usepackage{multido}
\usepackage{pst-fun}
\newcommand{\e}{\mathrm{e}}
\usepackage{amsmath}
\usepackage{amssymb}
\usepackage{etoolbox}
\newbool{dys}
\setbool{dys}{false}          
\ifbool{dys}{
% POLICE DYS
% ===== VARIABLE =====
\def\FontChoisie{Fira} % valeurs possibles : Fira, lmodern, tgheros
\newcommand{\choiceFontsDys}[1]{
\ifstrequal{#1}{Fira}{%
  % Fira Sans + Fira Math
  % Description : Fira Sans est une police moderne, et Fira Math est une version mathématique compatible qui maintient le style sans-serif.
  % Utilisation : Fonctionne bien avec XeLaTeX et LuaLaTeX.
  \usepackage{fontspec}
  \setmainfont{Fira Sans}
  \setsansfont{Fira Sans}
  \usepackage{unicode-math}
  \setmathfont{Fira Math}
}{}
\ifstrequal{#1}{lmodern}{%
  % Latin Modern Sans
  % Description : Une version sans-serif de la célèbre police Latin Modern, qui est compatible avec mathastext.
  % Utilisation : Fonctionne avec pdfLaTeX, XeLaTeX et LuaLaTeX.
  \usepackage{lmodern}
  \renewcommand{\familydefault}{\sfdefault}
  \usepackage{mathastext}
}{}
\ifstrequal{#1}{tgheros}{%
  % TeX Gyre Heros + mathastext
  % Description : TeX Gyre Heros est une alternative à Helvetica, et le package mathastext permet d'utiliser la même police pour les mathématiques.
  % Utilisation : Fonctionne avec pdfLaTeX, XeLaTeX, ou LuaLaTeX.
  \usepackage{tgheros}
  \renewcommand{\familydefault}{\sfdefault}
  \usepackage{mathastext}
}{}
}
% Fira ou lmodern ou tgheros
\choiceFontsDys{\FontChoisie}

%\usepackage{unicode-math}
%\usepackage{fontspec}
%\setmainfont{TeX Gyre Schola}
%\setsansfont{TeX Gyre Schola}
%\setmathfont{TeX Gyre Schola Math}
%\setmainfont{OpenDyslexic}[Scale=1.0]
%\setsansfont{OpenDyslexic}[Scale=1.0]
%\setmathfont{OpenDyslexic}[Scale=1.0,range=up/{Latin,latin,num}]
\usepackage[fontsize=14]{scrextend}
\usepackage{setspace}
\setstretch{1.7}
% Une valeur d'environ 1.2 à 1.7 est souvent recommandée. Cela permet d'augmenter l'espace entre les lignes tout en offrant un léger espacement entre les mots.
\setlength{\spaceskip}{1.2em}  
% Une valeur d'environ 1.2em à 1.5em est couramment conseillée. Cela crée un espace plus ample entre les mots, ce qui peut aider à réduire la fatigue visuelle et à améliorer la fluidité de la lecture.
}{
% POLICE STANDARD
\usepackage[T1]{fontenc} 
\usepackage[scaled=1]{helvet}
\usepackage[fontsize=13]{scrextend}
}
\begin{document}
\pagestyle{empty}
\begin{center}{\textbf{{\Huge \center Contrôle chapitres 7 et 8}\\
On prendra soin de rédiger correctement. Calculatrice autorisée.}.}\end{center}
\subsection*{Exercice 1 }
\bigskip

Calculer les longueurs des segments $[OQ]$ et $[PQ]$. Justifier.\\\begin{minipage}{0.5\linewidth}
    \begin{tikzpicture}[baseline,baseline=(current bounding box.north),scale = 0.5]

    \tikzset{
      point/.style={
        thick,
        draw,
        cross out,
        inner sep=0pt,
        minimum width=5pt,
        minimum height=5pt,
      },
    }
    \clip (-5.425584151765601,-5.8307452837322575) rectangle (1.0843027789744277,3.2682301564275855);
    	\draw[color={black}] (0,0)--(-4.133172,-4.75467)--(-4.586307,2.230648)--cycle;
	\draw[color={black},line width = 4] (-0.6560590289905075,-0.7547095802227717) arc (-131:-205.94:1) ;
	\draw[color={black}] (-3.5427187565487404,-4.075431733202968) arc (49:93.71000000000001:0.9) ;
	\draw[color={black}] (-3.411506950750639,-3.9244898171584133) arc (49:93.71000000000001:1.1) ;
	\draw[color={black}] (-4.528046659450559,1.33253540016608) arc (-86.29:-25.940000000000005:0.9) ;
	\draw[color={black}] (-4.5215733024536675,1.2327451418780548) arc (-86.29:-25.940000000000005:1) ;
	\draw[color={black}] (-4.515099945456776,1.1329548835900294) arc (-86.29:-25.940000000000005:1.1) ;
	\draw[opacity=1] (-1.6892311512087128,-2.705364692196985) node[anchor = center, rotate=48.99999999999997] {\normalsize \color{black}$6{,}3\text{ cm}$};
	\draw[opacity=1] (-2.0744624828036122,1.5649617910480342) node[anchor = center, rotate=-25.936955432503858] {\normalsize \color{black}$5{,}1\text{ cm}$};
	\draw[opacity=1] (-4.858690668971515,-1.294378100307033) node[anchor = center, rotate=-86.28844439152228] {\normalsize \color{black}$7\text{ cm}$};
	\draw [color={black}] (0.484303,0.140191) node[anchor = center,scale=1, rotate = 0] {H};
	\draw [color={black}] (-4.282403,-5.230745) node[anchor = center,scale=1, rotate = 0] {I};
	\draw [color={black}] (-4.825584,2.66823) node[anchor = center,scale=1, rotate = 0] {G};

\end{tikzpicture}\\    \end{minipage}
    \begin{minipage}{0.5\linewidth}
    \begin{tikzpicture}[baseline,baseline=(current bounding box.north),scale = 0.5]

    \tikzset{
      point/.style={
        thick,
        draw,
        cross out,
        inner sep=0pt,
        minimum width=5pt,
        minimum height=5pt,
      },
    }
    \clip (-6.121023983721285,-9.045586017185652) rectangle (5.412267502706067,6.850919763179667);
    	\draw[color={black}] (4.350658,-4.558399)--(-2.876405,5.762916)--(-5.23804,-8.036456)--cycle;
	\draw[color={black},line width = 4] (3.777081229755626,-3.7392473368136585) arc (125:199.94:1) ;
	\draw[color={black}] (-2.3601866392005637,5.025679537078555) arc (-55:-99.71000000000001:0.9) ;
	\draw[color={black}] (-2.2454713519303544,4.861849128220757) arc (-55:-99.71000000000001:1.1) ;
	\draw[color={black}] (-5.0862205398579015,-7.149353647681883) arc (80.29:19.940000000000005:0.9) ;
	\draw[color={black}] (-5.069351722240028,-7.050786700929055) arc (80.29:19.940000000000005:1) ;
	\draw[color={black}] (-5.052482904622154,-6.952219754176227) arc (80.29:19.940000000000005:1.1) ;
	\draw[opacity=1] (1.1467021392395795,0.8890467160935217) node[anchor = center, rotate=-55.00000000000002] {\normalsize \color{black}$18{,}9\text{ cm}$};
	\draw [color={black}] (4.812268,-4.746254) node[anchor = center,scale=1, rotate = 0] {O};
	\draw [color={black}] (-2.974841,6.25092) node[anchor = center,scale=1, rotate = 0] {P};
	\draw [color={black}] (-5.521024,-8.445586) node[anchor = center,scale=1, rotate = 0] {Q};

\end{tikzpicture}\\    \end{minipage}
 




\subsection*{Exercice 2 - D'après Métropole 2025}

\begin{minipage}{0.45\linewidth}
\begin{itemize}[label=$\bullet~$]
\item ABC un triangle rectangle en B ;
\item les points B, E et C sont alignés ainsi que
les points A, D, F et C ;
\item les droites (BD) et (EF) sont parallèles :
\item AB = 10 cm, BC = 7,5 cm, BE = 3 cm,

BD = 6 cm et CF = 2,7 cm.
\end{itemize}
\end{minipage}
\begin{minipage}{0.5\linewidth}
\psset{unit=0.9cm}
\begin{pspicture}(-0.5,-0.5)(8.5,6.5)
%\psgrid
\pspolygon(0,0)(7.8,0)(0,5.8)
\psline(2.78,3.75)
\psline(0,2.4)(1.63,4.6)
\uput[d](7.8,0){A}\uput[dl](0,0){B}\uput[ul](0,5.8){C}\uput[ur](2.78,3.75){D}
\uput[l](0,2.4){E}\uput[ur](1.63,4.6){F}
\psframe(0.3,0.3)
\end{pspicture}
\end{minipage}


\medskip

\begin{enumerate}
\item
	\begin{enumerate}
		\item Montrer que CE $= 4 ,5$ cm.
		\item Démontrer que la longueur EF est égale à $3,6$ cm.
	\end{enumerate}
\item Démontrer que le triangle CEF est rectangle en F.
\item  Les triangles ABC et CEF sont-ils semblables ? 
\end{enumerate}
 
\subsection*{Exercice 3 - CRPE 2025}
 
Un musée propose trois formules de visites guidées pour des classes durant l'année scolaire. 

\medskip
\fbox{
\begin{minipage}{\linewidth}
Formule A : 45~\euro{} par visite de classe. 

Formule B : Abonnement annuel de 90~\euro{} par école auquel s'ajoute un montant de 25~\euro{}~50 par visite de classe de l'école. 

Formule C : Abonnement annuel d'un montant de 300~\euro{} qui permet autant de visites que le souhaite l'école.	
\end{minipage}
}

\medskip
\begin{enumerate}
	\item Une école est composée de quatre classes.
	\begin{enumerate}
		\item Si chaque classe effectue une visite, quelle formule est la plus avantageuse ? 
		\item Si chaque classe effectue deux visites, quelle formule est la plus avantageuse ?
	\end{enumerate}
	\item On note $x$ le nombre de visites effectuées par une classe. Pour cette classe, on note $f(x)$ le prix avec la formule $A$, $g(x)$ le prix avec la formule B, et $h(x)$ le prix avec la formule $C$. Donner une expression littérale des fonctions $f$, $g$ et $h$ en fonction de $x$.
	\item Sur la représentation ci-dessous, associer à chacune des fonctions $f$, $g$ et $h$ la représentation qui lui convient.  
	
	\begin{tikzpicture}
\begin{axis}[
    axis lines = left,
    xmin=0, xmax=12,
    ymin=0, ymax=320,
    xtick={0,2,...,12},
    ytick={0,50,...,300},
    grid=major,
    width=16cm,
    height=15cm,
    legend pos=north west
    ]

\addplot [
    domain=0:12, 
    samples=100, 
    color=blue,    
    mark=o,
    mark size = 4pt,
    mark repeat = 7 pt,
    thick
    ]
    {45*x};
    \addlegendentry{$(d_1)$}

\addplot [
    domain=0:12, 
    samples=100, 
    thick, 
    mark =square*, mark repeat = 7, mark size = 2 pt,
    ]
    {90+25.5*x};
    \addlegendentry{$(d_2)$}

\addplot [
    domain=0:12, 
    samples=100, 
    color=red,
    thick
    ]
    {300};
    \addlegendentry{$(d_3)$}

\end{axis}
\end{tikzpicture}
\item Déterminer graphiquement le nombre de visites de classe à partir duquel la formule B est plus économique que la formule A.
	\item À partir de combien de visites de classe la formule C est-elle plus économique que la formule~B ? Vérifier par le calcul.

	\item L'équipe pédagogique de l'école de quatre classes décide d'organiser deux visites par classe. Elle choisit la formule la plus avantageuse et bénéficie d'une subvention de 15~\% du prix à payer de la part de la mairie. Quel montant l'école doit-elle prévoir pour payer le reste ?
\end{enumerate}



\end{document}